% !PRINTED 557
\weberpage{557}

\S.\,154.

\begin{center}
\gesperrt{The Group of the Ideal Classes.}
\end{center}

The ideal classes can be compounded on the basis of the following theorem:

1.~\gesperrt{If $\varphi$, $\psi$, $\varphi_1$, $\psi_1$ are functionals, and $\varphi$ is equivalent to $\varphi_1$, $\psi$ equivalent to $\psi_1$, then $\varphi\psi$ is also equivalent to $\varphi_1 \psi_1$.}

The correctness of this follows immediately from the definition. For if $\varphi:\varphi_1$ and $\psi:\psi_1$ are associated with numbers, then $\varphi\psi:\varphi_1\psi_1$ is also associated with a number.

In the same way the converse theorem also follows:

2.~\gesperrt{If $\varphi$ is equivalent to $\varphi_1$ and $\varphi\psi$ equivalent to $\varphi_1\psi_1$, then $\psi$ is also equivalent to $\psi_1$.}

If, then, we consider two ideal classes $A$, $B$, which may also be identical, and form the product $\varphi\psi$ of any functional $\varphi$ from $A$ and a functional $\psi$ from $B$, the class $C$ in which the product $\varphi\psi$ occurs is independent of the choice of $\varphi$ and $\psi$, and the class $C$ is completely determined by the two classes $A$, $B$. \gesperrt{We call $C$ compounded from $A$ and $B$}, and write symbolically

\weq{1}{C = AB = BA.}

The class $C$ contains all products of an element of $A$ with an element of $B$, but it may also contain still other functionals.

Since this composition is derived from true multiplication, the laws of multiplication also hold for this composition, namely the commutative and the associative law.

It follows further from Theorem 2. that if $AB = AB_1$, then $B = B_1$ must also hold, and consequently the ideal classes, under this composition, generate a \gesperrt{finite Abelian group of degree $h$}, to which we can apply all the theorems that we have come to know about such groups in the second chapter of this Band.

The unit of this group is the principal class $O$ already defined in \S.\,152, which indeed, as we have seen, has the represen-
%CONT>>
